\section{Memory encryption / integrity / replay protection}

\begin{frame}[t,shrink=6]{Memory encryption / integrity / replay protection：方向开场}
\DeckKicker{方向开场}
\SlideMeta{证据等级}{方向综述}{来源状态：公开材料综合}
{\large\bfseries\color{Ink} Memory encryption / integrity / replay protection 的核心是先界定保护对象、管理边界和证据强度，再用三篇主讲材料串起技术演进。\par}\vspace{1.8mm}
\begin{columns}[T,totalwidth=\textwidth]
  \begin{column}{0.45\textwidth}
    \fcolorbox{Rule}{Paper}{%
  \begin{minipage}[t]{0.97\linewidth}
    \vspace{1.1mm}
    \hspace{1.4mm}\begin{minipage}{0.92\linewidth}
      \PanelTitle{内容说明}
      \scriptsize \begin{itemize}
  \item 避免把 access control、ownership、encryption、integrity、freshness/replay 混写。
  \item 固定三篇主讲材料；`riscv\_ap\_tee\_2024` 作为 AP-TEE/CoVE lifecycle auxiliary，不计入三篇主讲。
  \item 先讲方向痛点和设计轴，再逐篇说明贡献、限制、证据强度和可落地条件。
  \item 对规范、Survey、draft、vendor evidence 保持显式标注，避免把背景材料写成一手系统证明。
\end{itemize}
    \end{minipage}
    \vspace{1mm}
  \end{minipage}}
  \end{column}
  \begin{column}{0.49\textwidth}
    \fcolorbox{Rule}{Panel}{%
  \begin{minipage}[t]{0.97\linewidth}
    \vspace{1.1mm}
    \hspace{1.4mm}\begin{minipage}{0.92\linewidth}
      \PanelTitle{三篇主讲材料的分工}
      \scriptsize {\tiny
\begin{tabularx}{\linewidth}{@{}YYY@{}}
\toprule
\textbf{论文} & \textbf{定位} & \textbf{证据边界} \\
\midrule
Memory Encryption: A Survey of Existing Techniques & 开创/基础入口 & E2 / Background substrate \\
Bonsai Merkle Trees for memory integrity and freshness & 代表性改进 & E1 / Peer-reviewed mechanism anchor \\
AMD SEV-SNP industry memory-integrity evidence & 当前 SOTA/补充边界 & E4 / Industry evidence \\
\bottomrule
\end{tabularx}
}
    \end{minipage}
    \vspace{1mm}
  \end{minipage}}
  \end{column}
\end{columns}
\vspace{1mm}
{\tiny\textcolor{Muted}{引用依据：本页综合三篇主讲材料的研究定位、证据等级和公开来源状态。}}
\end{frame}


\begin{frame}[t,shrink=6]{Memory Encryption: A Survey of Existing Techniques：内容摘要}
\DeckKicker{内容摘要}
\SlideMeta{证据等级}{E2 / Background substrate}{来源状态：本地 PDF 已核验}
{\large\bfseries\color{Ink} 综述 memory encryption 的威胁模型、技术路线和性能/安全权衡\par}\vspace{1.8mm}
\begin{columns}[T,totalwidth=\textwidth]
  \begin{column}{0.45\textwidth}
    \fcolorbox{Rule}{Paper}{%
  \begin{minipage}[t]{0.97\linewidth}
    \vspace{1.1mm}
    \hspace{1.4mm}\begin{minipage}{0.92\linewidth}
      \PanelTitle{内容说明}
      \scriptsize \begin{itemize}
  \item 综述 memory encryption 的威胁模型、技术路线和性能/安全权衡
  \item Henson 2014 is the background starting point for secure memory encryption, integrity, and r。
\end{itemize}
    \end{minipage}
    \vspace{1mm}
  \end{minipage}}
  \end{column}
  \begin{column}{0.49\textwidth}
    \fcolorbox{Rule}{Panel}{%
  \begin{minipage}[t]{0.97\linewidth}
    \vspace{1.1mm}
    \hspace{1.4mm}\begin{minipage}{0.92\linewidth}
      \PanelTitle{证据定位}
      \scriptsize {\tiny
\begin{tabularx}{\linewidth}{@{}YY@{}}
\toprule
\textbf{字段} & \textbf{内容} \\
\midrule
论文定位 & 开创/基础入口 \\
材料类型 & SoK/综述 \\
证据等级 & E2 / Background substrate \\
来源状态 & 本地 PDF 已核验 \\
\bottomrule
\end{tabularx}
}
    \end{minipage}
    \vspace{1mm}
  \end{minipage}}
  \end{column}
\end{columns}
\vspace{1mm}
{\tiny\textcolor{Muted}{引用依据：引用依据为论文原文、官方规范或公开材料；本页结论按 E2 / Background substrate 的证据等级使用，不外推到未验证平台。}}
\end{frame}


\begin{frame}[t,shrink=6]{Memory Encryption: A Survey of Existing Techniques：研究背景}
\DeckKicker{研究背景}
\SlideMeta{证据等级}{E2 / Background substrate}{来源状态：本地 PDF 已核验}
{\large\bfseries\color{Ink} Full disk encryption 不能保护运行时 DRAM、bus 和 DMA 路径上的明文\par}\vspace{1.8mm}
\begin{columns}[T,totalwidth=\textwidth]
  \begin{column}{0.45\textwidth}
    \fcolorbox{Rule}{Paper}{%
  \begin{minipage}[t]{0.97\linewidth}
    \vspace{1.1mm}
    \hspace{1.4mm}\begin{minipage}{0.92\linewidth}
      \PanelTitle{内容说明}
      \scriptsize \begin{itemize}
  \item Full disk encryption 不能保护运行时 DRAM、bus 和 DMA 路径上的明文
  \item access control 与 encryption/integrity/replay 是不同问题
\end{itemize}
    \end{minipage}
    \vspace{1mm}
  \end{minipage}}
  \end{column}
  \begin{column}{0.49\textwidth}
    \fcolorbox{Rule}{Panel}{%
  \begin{minipage}[t]{0.97\linewidth}
    \vspace{1.1mm}
    \hspace{1.4mm}\begin{minipage}{0.92\linewidth}
      \PanelTitle{问题边界}
      \scriptsize {\tiny
\begin{tabularx}{\linewidth}{@{}YY@{}}
\toprule
\textbf{维度} & \textbf{说明} \\
\midrule
保护对象 & Memory encryption / integrity / replay protection \\
传统瓶颈 & 避免把 access control、ownership、encryption、integrity、freshness/replay 混写。 \\
本文入口 & Memory Encryption: A Survey of Existing Techniques \\
证据限制 & E2 / Background substrate \\
\bottomrule
\end{tabularx}
}
    \end{minipage}
    \vspace{1mm}
  \end{minipage}}
  \end{column}
\end{columns}
\vspace{1mm}
{\tiny\textcolor{Muted}{引用依据：引用依据为论文原文、官方规范或公开材料；本页结论按 E2 / Background substrate 的证据等级使用，不外推到未验证平台。}}
\end{frame}


\begin{frame}[t,shrink=6]{Memory Encryption: A Survey of Existing Techniques：解决方案}
\DeckKicker{解决方案}
\SlideMeta{证据等级}{E2 / Background substrate}{来源状态：本地 PDF 已核验}
{\large\bfseries\color{Ink} 按硬件增强、OS 辅助、专用设备与 cryptographic memory controller 路线组织 taxonomy，并解释 encry。\par}\vspace{1.8mm}
\begin{columns}[T,totalwidth=\textwidth]
  \begin{column}{0.45\textwidth}
    \fcolorbox{Rule}{Paper}{%
  \begin{minipage}[t]{0.97\linewidth}
    \vspace{1.1mm}
    \hspace{1.4mm}\begin{minipage}{0.92\linewidth}
      \PanelTitle{内容说明}
      \scriptsize \begin{itemize}
  \item 按硬件增强、OS 辅助、专用设备与 cryptographic memory controller 路线组织 taxonomy，并解释 encryption、authenticati。
\end{itemize}
    \end{minipage}
    \vspace{1mm}
  \end{minipage}}
  \end{column}
  \begin{column}{0.49\textwidth}
    \fcolorbox{Rule}{Panel}{%
  \begin{minipage}[t]{0.97\linewidth}
    \vspace{1.1mm}
    \hspace{1.4mm}\begin{minipage}{0.92\linewidth}
      \PanelTitle{核心设计拆解}
      \scriptsize \begin{itemize}
  \item 01. Off-chip memory confidentiality threat model
  \item 02. Memory authentication and integrity taxonomy
  \item 03. Counter/freshness/replay-protection tradeoff
\end{itemize}
    \end{minipage}
    \vspace{1mm}
  \end{minipage}}
  \end{column}
\end{columns}
\vspace{1mm}
{\tiny\textcolor{Muted}{引用依据：引用依据为论文原文、官方规范或公开材料；本页结论按 E2 / Background substrate 的证据等级使用，不外推到未验证平台。}}
\end{frame}


\begin{frame}[t,shrink=6]{Memory Encryption: A Survey of Existing Techniques：实验结果}
\DeckKicker{实验结果}
\SlideMeta{证据等级}{E2 / Background substrate}{来源状态：本地 PDF 已核验}
{\large\bfseries\color{Ink} 规范/Survey，无新实验；Memory Encryption: A Survey of Existing Techniques 只支撑术语、taxonomy、接口语义或证据边界。\par}\vspace{1.8mm}
\begin{columns}[T,totalwidth=\textwidth]
  \begin{column}{0.45\textwidth}
    \fcolorbox{Rule}{Paper}{%
  \begin{minipage}[t]{0.97\linewidth}
    \vspace{1.1mm}
    \hspace{1.4mm}\begin{minipage}{0.92\linewidth}
      \PanelTitle{内容说明}
      \scriptsize \begin{itemize}
  \item 本页不展示独立系统实验，也不把二手归纳写成一手结果。
  \item 可支撑：Henson 2014 is the background starting point for secure memory encryption, integrity。
  \item 证据等级：E2 / Background substrate；来源状态：本地 PDF 已核验。
  \item 涉及具体平台、协议状态或数值时，转到原始论文、官方规范或厂商材料核验。
\end{itemize}
    \end{minipage}
    \vspace{1mm}
  \end{minipage}}
  \end{column}
  \begin{column}{0.49\textwidth}
    \fcolorbox{Rule}{Panel}{%
  \begin{minipage}[t]{0.97\linewidth}
    \vspace{1.1mm}
    \hspace{1.4mm}\begin{minipage}{0.92\linewidth}
      \PanelTitle{实验/证据边界}
      \scriptsize {\tiny
\begin{tabularx}{\linewidth}{@{}YYY@{}}
\toprule
\textbf{对象} & \textbf{可支撑} & \textbf{边界} \\
\midrule
数据/来源 & Memory Encryption: A Survey of Existing Techniques & 本地 PDF 已核验 \\
Baseline/对照 & 以论文或规范原文为准 & 不跨材料外推 \\
指标/对象 & 只使用当前来源可证明的信息 & E2 / Background substrate \\
使用边界 & 可进入本方向演进图 & 不能替代更强证据 \\
\bottomrule
\end{tabularx}
}
    \end{minipage}
    \vspace{1mm}
  \end{minipage}}
  \end{column}
\end{columns}
\vspace{1mm}
{\tiny\textcolor{Muted}{引用依据：引用依据为论文原文、官方规范或公开材料；本页结论按 E2 / Background substrate 的证据等级使用，不外推到未验证平台。}}
\end{frame}


\begin{frame}[t,shrink=6]{Memory Encryption: A Survey of Existing Techniques：文章评价}
\DeckKicker{文章评价}
\SlideMeta{证据等级}{E2 / Background substrate}{来源状态：本地 PDF 已核验}
{\large\bfseries\color{Ink} 评价：Memory Encryption: A Survey of Existing Techniques 适合做 taxonomy/规范入口；规范/Survey，无新实验，不能替代一手系统证据。\par}\vspace{1.8mm}
\begin{columns}[T,totalwidth=\textwidth]
  \begin{column}{0.45\textwidth}
    \fcolorbox{Rule}{Paper}{%
  \begin{minipage}[t]{0.97\linewidth}
    \vspace{1.1mm}
    \hspace{1.4mm}\begin{minipage}{0.92\linewidth}
      \PanelTitle{内容说明}
      \scriptsize \begin{itemize}
  \item 设计优势：最适合给本方向提供概念边界，防止把 PMP/GPT/IOPMP 写成 encryption。
  \item 局限性：2014 时间截面早于 SEV-SNP、TDX、Arm CCA、RISC-V CoVE 和 CXL/PCIe IDE。
  \item 商业化潜力：可作为架构设计 checklist；风险是单独引用会遗漏现代 confidential VM ownership metadata。
  \item 规范/Survey，无新实验；具体平台结论转到一手材料核验。
\end{itemize}
    \end{minipage}
    \vspace{1mm}
  \end{minipage}}
  \end{column}
  \begin{column}{0.49\textwidth}
    \fcolorbox{Rule}{Panel}{%
  \begin{minipage}[t]{0.97\linewidth}
    \vspace{1.1mm}
    \hspace{1.4mm}\begin{minipage}{0.92\linewidth}
      \PanelTitle{文章评价}
      \scriptsize {\tiny
\begin{tabularx}{\linewidth}{@{}YY@{}}
\toprule
\textbf{维度} & \textbf{判断} \\
\midrule
设计优势 & 最适合给本方向提供概念边界，防止把 PMP/GPT/IOPMP 写成 encryption。 \\
主要局限 & 2014 时间截面早于 SEV-SNP、TDX、Arm CCA、RISC-V CoVE 和 CXL/PCIe IDE。 \\
商业落地 & 可作为架构设计 checklist；风险是单独引用会遗漏现代 confidential VM ownership metadata。 \\
讲稿定位 & 开创/基础入口; 证据强度 E2 \\
\bottomrule
\end{tabularx}
}
    \end{minipage}
    \vspace{1mm}
  \end{minipage}}
  \end{column}
\end{columns}
\vspace{1mm}
{\tiny\textcolor{Muted}{引用依据：引用依据为论文原文、官方规范或公开材料；本页结论按 E2 / Background substrate 的证据等级使用，不外推到未验证平台。}}
\end{frame}


\begin{frame}[t,shrink=6]{Bonsai Merkle Trees for memory integrity and freshness：内容摘要}
\DeckKicker{内容摘要}
\SlideMeta{证据等级}{E1 / Peer-reviewed mechanism anchor}{来源状态：本地 PDF 已核验}
{\large\bfseries\color{Ink} Bonsai Merkle Trees use counter/tree metadata to detect tamper and replay, anchoring memory freshness evidence.\par}\vspace{1.8mm}
\begin{columns}[T,totalwidth=\textwidth]
  \begin{column}{0.45\textwidth}
    \fcolorbox{Rule}{Paper}{%
  \begin{minipage}[t]{0.97\linewidth}
    \vspace{1.1mm}
    \hspace{1.4mm}\begin{minipage}{0.92\linewidth}
      \PanelTitle{内容说明}
      \scriptsize \begin{itemize}
  \item Counter/tree metadata detects tamper and replay
  \item Provides peer-reviewed evidence for memory integrity and freshness
  \item Peer-reviewed Bonsai Merkle Tree mechanism anchor for memory integrity, freshness, and repl。
\end{itemize}
    \end{minipage}
    \vspace{1mm}
  \end{minipage}}
  \end{column}
  \begin{column}{0.49\textwidth}
    \fcolorbox{Rule}{Panel}{%
  \begin{minipage}[t]{0.97\linewidth}
    \vspace{1.1mm}
    \hspace{1.4mm}\begin{minipage}{0.92\linewidth}
      \PanelTitle{证据定位}
      \scriptsize {\tiny
\begin{tabularx}{\linewidth}{@{}YY@{}}
\toprule
\textbf{字段} & \textbf{内容} \\
\midrule
论文定位 & 代表性改进 \\
材料类型 & 系统论文 \\
证据等级 & E1 / Peer-reviewed mechanism anchor \\
来源状态 & 本地 PDF 已核验 \\
\bottomrule
\end{tabularx}
}
    \end{minipage}
    \vspace{1mm}
  \end{minipage}}
  \end{column}
\end{columns}
\vspace{1mm}
{\tiny\textcolor{Muted}{引用依据：引用依据为论文原文、官方规范或公开材料；本页结论按 E1 / Peer-reviewed mechanism anchor 的证据等级使用，不外推到未验证平台。}}
\end{frame}


\begin{frame}[t,shrink=6]{Bonsai Merkle Trees for memory integrity and freshness：研究背景}
\DeckKicker{研究背景}
\SlideMeta{证据等级}{E1 / Peer-reviewed mechanism anchor}{来源状态：本地 PDF 已核验}
{\large\bfseries\color{Ink} Encryption alone cannot detect stale or modified memory values returned by an untrusted memory system.\par}\vspace{1.8mm}
\begin{columns}[T,totalwidth=\textwidth]
  \begin{column}{0.45\textwidth}
    \fcolorbox{Rule}{Paper}{%
  \begin{minipage}[t]{0.97\linewidth}
    \vspace{1.1mm}
    \hspace{1.4mm}\begin{minipage}{0.92\linewidth}
      \PanelTitle{内容说明}
      \scriptsize \begin{itemize}
  \item Encryption alone cannot detect stale or modified memory values
  \item Freshness requires integrity trees, counters, or authenticated roots
\end{itemize}
    \end{minipage}
    \vspace{1mm}
  \end{minipage}}
  \end{column}
  \begin{column}{0.49\textwidth}
    \fcolorbox{Rule}{Panel}{%
  \begin{minipage}[t]{0.97\linewidth}
    \vspace{1.1mm}
    \hspace{1.4mm}\begin{minipage}{0.92\linewidth}
      \PanelTitle{问题边界}
      \scriptsize {\tiny
\begin{tabularx}{\linewidth}{@{}YY@{}}
\toprule
\textbf{维度} & \textbf{说明} \\
\midrule
保护对象 & Memory encryption / integrity / replay protection \\
传统瓶颈 & 避免把 access control、ownership、encryption、integrity、freshness/replay 混写。 \\
本文入口 & Bonsai Merkle Trees for memory integrity and freshness \\
证据限制 & E1 / Peer-reviewed mechanism anchor \\
\bottomrule
\end{tabularx}
}
    \end{minipage}
    \vspace{1mm}
  \end{minipage}}
  \end{column}
\end{columns}
\vspace{1mm}
{\tiny\textcolor{Muted}{引用依据：引用依据为论文原文、官方规范或公开材料；本页结论按 E1 / Peer-reviewed mechanism anchor 的证据等级使用，不外推到未验证平台。}}
\end{frame}


\begin{frame}[t,shrink=6]{Bonsai Merkle Trees for memory integrity and freshness：解决方案}
\DeckKicker{解决方案}
\SlideMeta{证据等级}{E1 / Peer-reviewed mechanism anchor}{来源状态：本地 PDF 已核验}
{\large\bfseries\color{Ink} Address-independent seed encryption and Bonsai Merkle Trees reduce secure-processor integrity-tree overhead.\par}\vspace{1.8mm}
\begin{columns}[T,totalwidth=\textwidth]
  \begin{column}{0.45\textwidth}
    \fcolorbox{Rule}{Paper}{%
  \begin{minipage}[t]{0.97\linewidth}
    \vspace{1.1mm}
    \hspace{1.4mm}\begin{minipage}{0.92\linewidth}
      \PanelTitle{内容说明}
      \scriptsize \begin{itemize}
  \item Address-independent seed encryption reduces OS friction
  \item Bonsai Merkle Trees authenticate counter metadata and root state
\end{itemize}
    \end{minipage}
    \vspace{1mm}
  \end{minipage}}
  \end{column}
  \begin{column}{0.49\textwidth}
    \fcolorbox{Rule}{Panel}{%
  \begin{minipage}[t]{0.97\linewidth}
    \vspace{1.1mm}
    \hspace{1.4mm}\begin{minipage}{0.92\linewidth}
      \PanelTitle{核心设计拆解}
      \scriptsize \begin{itemize}
  \item 01. AISE counter-mode seed
  \item 02. BMT over encryption counters
  \item 03. root MAC for replay/tamper detection
\end{itemize}
    \end{minipage}
    \vspace{1mm}
  \end{minipage}}
  \end{column}
\end{columns}
\vspace{1mm}
{\tiny\textcolor{Muted}{引用依据：引用依据为论文原文、官方规范或公开材料；本页结论按 E1 / Peer-reviewed mechanism anchor 的证据等级使用，不外推到未验证平台。}}
\end{frame}


\begin{frame}[t,shrink=6]{Bonsai Merkle Trees for memory integrity and freshness：实验结果}
\DeckKicker{实验结果}
\SlideMeta{证据等级}{E1 / Peer-reviewed mechanism anchor}{来源状态：本地 PDF 已核验}
{\large\bfseries\color{Ink} MICRO 2007 peer-reviewed paper reports secure-processor simulation results and reduced memory-integrity overhea。\par}\vspace{1.8mm}
\begin{columns}[T,totalwidth=\textwidth]
  \begin{column}{0.45\textwidth}
    \fcolorbox{Rule}{Paper}{%
  \begin{minipage}[t]{0.97\linewidth}
    \vspace{1.1mm}
    \hspace{1.4mm}\begin{minipage}{0.92\linewidth}
      \PanelTitle{内容说明}
      \scriptsize \begin{itemize}
  \item MICRO 2007 peer-reviewed secure-processor simulation
  \item Historical performance numbers are lineage evidence, not direct CCA/SEV-SNP product compari。
\end{itemize}
    \end{minipage}
    \vspace{1mm}
  \end{minipage}}
  \end{column}
  \begin{column}{0.49\textwidth}
    \fcolorbox{Rule}{Panel}{%
  \begin{minipage}[t]{0.97\linewidth}
    \vspace{1.1mm}
    \hspace{1.4mm}\begin{minipage}{0.92\linewidth}
      \PanelTitle{实验/证据边界}
      \scriptsize {\tiny
\begin{tabularx}{\linewidth}{@{}YYY@{}}
\toprule
\textbf{对象} & \textbf{可支撑} & \textbf{边界} \\
\midrule
数据/来源 & Bonsai Merkle Trees for memory integrity and freshness & 本地 PDF 已核验 \\
Baseline/对照 & 以论文或规范原文为准 & 不跨材料外推 \\
指标/对象 & 只使用当前来源可证明的信息 & E1 / Peer-reviewed mechanism anchor \\
使用边界 & 可进入本方向演进图 & 不能替代更强证据 \\
\bottomrule
\end{tabularx}
}
    \end{minipage}
    \vspace{1mm}
  \end{minipage}}
  \end{column}
\end{columns}
\vspace{1mm}
{\tiny\textcolor{Muted}{引用依据：引用依据为论文原文、官方规范或公开材料；本页结论按 E1 / Peer-reviewed mechanism anchor 的证据等级使用，不外推到未验证平台。}}
\end{frame}


\begin{frame}[t,shrink=6]{Bonsai Merkle Trees for memory integrity and freshness：文章评价}
\DeckKicker{文章评价}
\SlideMeta{证据等级}{E1 / Peer-reviewed mechanism anchor}{来源状态：本地 PDF 已核验}
{\large\bfseries\color{Ink} Moves the section anchor from vendor/spec contrast back to peer-reviewed memory-integrity and freshness mechani。\par}\vspace{1.8mm}
\begin{columns}[T,totalwidth=\textwidth]
  \begin{column}{0.45\textwidth}
    \fcolorbox{Rule}{Paper}{%
  \begin{minipage}[t]{0.97\linewidth}
    \vspace{1.1mm}
    \hspace{1.4mm}\begin{minipage}{0.92\linewidth}
      \PanelTitle{内容说明}
      \scriptsize \begin{itemize}
  \item 设计优势：Moves the section anchor from vendor/spec contrast back to peer-reviewed memory-integ。
  \item 局限性：Secure-processor lineage is not the same as a modern confidential-VM product; deploym。
  \item 商业化潜力：Useful for memory-controller integrity-tree design; productization depends on metadat。
  \item 证据边界：E1 / Peer-reviewed mechanism anchor；本地 PDF 已核验。
\end{itemize}
    \end{minipage}
    \vspace{1mm}
  \end{minipage}}
  \end{column}
  \begin{column}{0.49\textwidth}
    \fcolorbox{Rule}{Panel}{%
  \begin{minipage}[t]{0.97\linewidth}
    \vspace{1.1mm}
    \hspace{1.4mm}\begin{minipage}{0.92\linewidth}
      \PanelTitle{文章评价}
      \scriptsize {\tiny
\begin{tabularx}{\linewidth}{@{}YY@{}}
\toprule
\textbf{维度} & \textbf{判断} \\
\midrule
设计优势 & Moves the section anchor from vendor/spec contrast back to peer-reviewed memory。 \\
主要局限 & Secure-processor lineage is not the same as a modern confidential-VM product; d。 \\
商业落地 & Useful for memory-controller integrity-tree design; productization depends on m。 \\
讲稿定位 & 代表性改进; 证据强度 E1 \\
\bottomrule
\end{tabularx}
}
    \end{minipage}
    \vspace{1mm}
  \end{minipage}}
  \end{column}
\end{columns}
\vspace{1mm}
{\tiny\textcolor{Muted}{引用依据：引用依据为论文原文、官方规范或公开材料；本页结论按 E1 / Peer-reviewed mechanism anchor 的证据等级使用，不外推到未验证平台。}}
\end{frame}


\begin{frame}[t,shrink=6]{AMD SEV-SNP industry memory-integrity evidence：内容摘要}
\DeckKicker{内容摘要}
\SlideMeta{证据等级}{E4 / Industry evidence}{来源状态：本地 PDF 已核验}
{\large\bfseries\color{Ink} SEV-SNP 在 encrypted VM 基础上加入 secure nested paging 与完整性相关保护，减少 host-contro。\par}\vspace{1.8mm}
\begin{columns}[T,totalwidth=\textwidth]
  \begin{column}{0.45\textwidth}
    \fcolorbox{Rule}{Paper}{%
  \begin{minipage}[t]{0.97\linewidth}
    \vspace{1.1mm}
    \hspace{1.4mm}\begin{minipage}{0.92\linewidth}
      \PanelTitle{内容说明}
      \scriptsize \begin{itemize}
  \item SEV-SNP 在 encrypted VM 基础上加入 secure nested paging 与完整性相关保护，减少 host-controlled metadata 风险
  \item AMD SEV-SNP remains industry evidence for modern VM memory-integrity deployment, but not th。
\end{itemize}
    \end{minipage}
    \vspace{1mm}
  \end{minipage}}
  \end{column}
  \begin{column}{0.49\textwidth}
    \fcolorbox{Rule}{Panel}{%
  \begin{minipage}[t]{0.97\linewidth}
    \vspace{1.1mm}
    \hspace{1.4mm}\begin{minipage}{0.92\linewidth}
      \PanelTitle{证据定位}
      \scriptsize {\tiny
\begin{tabularx}{\linewidth}{@{}YY@{}}
\toprule
\textbf{字段} & \textbf{内容} \\
\midrule
论文定位 & 当前 SOTA/补充边界 \\
材料类型 & 厂商白皮书 \\
证据等级 & E4 / Industry evidence \\
来源状态 & 本地 PDF 已核验 \\
\bottomrule
\end{tabularx}
}
    \end{minipage}
    \vspace{1mm}
  \end{minipage}}
  \end{column}
\end{columns}
\vspace{1mm}
{\tiny\textcolor{Muted}{引用依据：引用依据为论文原文、官方规范或公开材料；本页结论按 E4 / Industry evidence 的证据等级使用，不外推到未验证平台。}}
\end{frame}


\begin{frame}[t,shrink=6]{AMD SEV-SNP industry memory-integrity evidence：研究背景}
\DeckKicker{研究背景}
\SlideMeta{证据等级}{E4 / Industry evidence}{来源状态：本地 PDF 已核验}
{\large\bfseries\color{Ink} 早期 encrypted VM 面临 nested page table、reverse map metadata、replay/alias 和。\par}\vspace{1.8mm}
\begin{columns}[T,totalwidth=\textwidth]
  \begin{column}{0.45\textwidth}
    \fcolorbox{Rule}{Paper}{%
  \begin{minipage}[t]{0.97\linewidth}
    \vspace{1.1mm}
    \hspace{1.4mm}\begin{minipage}{0.92\linewidth}
      \PanelTitle{内容说明}
      \scriptsize \begin{itemize}
  \item 早期 encrypted VM 面临 nested page table、reverse map metadata、replay/alias 和 malicious hypervis。
\end{itemize}
    \end{minipage}
    \vspace{1mm}
  \end{minipage}}
  \end{column}
  \begin{column}{0.49\textwidth}
    \fcolorbox{Rule}{Panel}{%
  \begin{minipage}[t]{0.97\linewidth}
    \vspace{1.1mm}
    \hspace{1.4mm}\begin{minipage}{0.92\linewidth}
      \PanelTitle{问题边界}
      \scriptsize {\tiny
\begin{tabularx}{\linewidth}{@{}YY@{}}
\toprule
\textbf{维度} & \textbf{说明} \\
\midrule
保护对象 & Memory encryption / integrity / replay protection \\
传统瓶颈 & 避免把 access control、ownership、encryption、integrity、freshness/replay 混写。 \\
本文入口 & AMD SEV-SNP industry memory-integrity evidence \\
证据限制 & E4 / Industry evidence \\
\bottomrule
\end{tabularx}
}
    \end{minipage}
    \vspace{1mm}
  \end{minipage}}
  \end{column}
\end{columns}
\vspace{1mm}
{\tiny\textcolor{Muted}{引用依据：引用依据为论文原文、官方规范或公开材料；本页结论按 E4 / Industry evidence 的证据等级使用，不外推到未验证平台。}}
\end{frame}


\begin{frame}[t,shrink=6]{AMD SEV-SNP industry memory-integrity evidence：解决方案}
\DeckKicker{解决方案}
\SlideMeta{证据等级}{E4 / Industry evidence}{来源状态：本地 PDF 已核验}
{\large\bfseries\color{Ink} 引入 secure nested paging、RMP 等机制，把 memory ownership/validation 与 encrypted。\par}\vspace{1.8mm}
\begin{columns}[T,totalwidth=\textwidth]
  \begin{column}{0.45\textwidth}
    \fcolorbox{Rule}{Paper}{%
  \begin{minipage}[t]{0.97\linewidth}
    \vspace{1.1mm}
    \hspace{1.4mm}\begin{minipage}{0.92\linewidth}
      \PanelTitle{内容说明}
      \scriptsize \begin{itemize}
  \item 引入 secure nested paging、RMP 等机制，把 memory ownership/validation 与 encrypted VM lifecycle 结合
\end{itemize}
    \end{minipage}
    \vspace{1mm}
  \end{minipage}}
  \end{column}
  \begin{column}{0.49\textwidth}
    \fcolorbox{Rule}{Panel}{%
  \begin{minipage}[t]{0.97\linewidth}
    \vspace{1.1mm}
    \hspace{1.4mm}\begin{minipage}{0.92\linewidth}
      \PanelTitle{核心设计拆解}
      \scriptsize \begin{itemize}
  \item 01. Secure Nested Paging
  \item 02. Reverse Map Table / memory ownership metadata
  \item 03. Integrity-oriented VM isolation controls
\end{itemize}
    \end{minipage}
    \vspace{1mm}
  \end{minipage}}
  \end{column}
\end{columns}
\vspace{1mm}
{\tiny\textcolor{Muted}{引用依据：引用依据为论文原文、官方规范或公开材料；本页结论按 E4 / Industry evidence 的证据等级使用，不外推到未验证平台。}}
\end{frame}


\begin{frame}[t,shrink=6]{AMD SEV-SNP industry memory-integrity evidence：实验结果}
\DeckKicker{实验结果}
\SlideMeta{证据等级}{E4 / Industry evidence}{来源状态：本地 PDF 已核验}
{\large\bfseries\color{Ink} 厂商技术材料，无独立学术实验\par}\vspace{1.8mm}
\begin{columns}[T,totalwidth=\textwidth]
  \begin{column}{0.45\textwidth}
    \fcolorbox{Rule}{Paper}{%
  \begin{minipage}[t]{0.97\linewidth}
    \vspace{1.1mm}
    \hspace{1.4mm}\begin{minipage}{0.92\linewidth}
      \PanelTitle{内容说明}
      \scriptsize \begin{itemize}
  \item 厂商技术材料，无独立学术实验
  \item 只能作为产业机制和部署 evidence
\end{itemize}
    \end{minipage}
    \vspace{1mm}
  \end{minipage}}
  \end{column}
  \begin{column}{0.49\textwidth}
    \fcolorbox{Rule}{Panel}{%
  \begin{minipage}[t]{0.97\linewidth}
    \vspace{1.1mm}
    \hspace{1.4mm}\begin{minipage}{0.92\linewidth}
      \PanelTitle{实验/证据边界}
      \scriptsize {\tiny
\begin{tabularx}{\linewidth}{@{}YYY@{}}
\toprule
\textbf{对象} & \textbf{可支撑} & \textbf{边界} \\
\midrule
数据/来源 & AMD SEV-SNP industry memory-integrity evidence & 本地 PDF 已核验 \\
Baseline/对照 & 以论文或规范原文为准 & 不跨材料外推 \\
指标/对象 & 只使用当前来源可证明的信息 & E4 / Industry evidence \\
使用边界 & 可进入本方向演进图 & 不能替代更强证据 \\
\bottomrule
\end{tabularx}
}
    \end{minipage}
    \vspace{1mm}
  \end{minipage}}
  \end{column}
\end{columns}
\vspace{1mm}
{\tiny\textcolor{Muted}{引用依据：引用依据为论文原文、官方规范或公开材料；本页结论按 E4 / Industry evidence 的证据等级使用，不外推到未验证平台。}}
\end{frame}


\begin{frame}[t,shrink=6]{AMD SEV-SNP industry memory-integrity evidence：文章评价}
\DeckKicker{文章评价}
\SlideMeta{证据等级}{E4 / Industry evidence}{来源状态：本地 PDF 已核验}
{\large\bfseries\color{Ink} 展示 memory encryption 从 confidentiality 走向 integrity/metadata protection 的。\par}\vspace{1.8mm}
\begin{columns}[T,totalwidth=\textwidth]
  \begin{column}{0.45\textwidth}
    \fcolorbox{Rule}{Paper}{%
  \begin{minipage}[t]{0.97\linewidth}
    \vspace{1.1mm}
    \hspace{1.4mm}\begin{minipage}{0.92\linewidth}
      \PanelTitle{内容说明}
      \scriptsize \begin{itemize}
  \item 设计优势：展示 memory encryption 从 confidentiality 走向 integrity/metadata protection 的商用方向。
  \item 局限性：Vendor whitepaper 不是 peer-reviewed proof，也不证明所有 SEV-SNP 部署抵抗实现攻击或侧信道。
  \item 商业化潜力：具有真实云平台相关性；风险在供应商绑定、固件/微码版本和 verifier policy 透明度。
  \item 证据边界：E4 / Industry evidence；本地 PDF 已核验。
\end{itemize}
    \end{minipage}
    \vspace{1mm}
  \end{minipage}}
  \end{column}
  \begin{column}{0.49\textwidth}
    \fcolorbox{Rule}{Panel}{%
  \begin{minipage}[t]{0.97\linewidth}
    \vspace{1.1mm}
    \hspace{1.4mm}\begin{minipage}{0.92\linewidth}
      \PanelTitle{文章评价}
      \scriptsize {\tiny
\begin{tabularx}{\linewidth}{@{}YY@{}}
\toprule
\textbf{维度} & \textbf{判断} \\
\midrule
设计优势 & 展示 memory encryption 从 confidentiality 走向 integrity/metadata protection 的商用方向。 \\
主要局限 & Vendor whitepaper 不是 peer-reviewed proof，也不证明所有 SEV-SNP 部署抵抗实现攻击或侧信道。 \\
商业落地 & 具有真实云平台相关性；风险在供应商绑定、固件/微码版本和 verifier policy 透明度。 \\
讲稿定位 & 当前 SOTA/补充边界; 证据强度 E4 \\
\bottomrule
\end{tabularx}
}
    \end{minipage}
    \vspace{1mm}
  \end{minipage}}
  \end{column}
\end{columns}
\vspace{1mm}
{\tiny\textcolor{Muted}{引用依据：引用依据为论文原文、官方规范或公开材料；本页结论按 E4 / Industry evidence 的证据等级使用，不外推到未验证平台。}}
\end{frame}


\begin{frame}[t,shrink=6]{Memory encryption / integrity / replay protection：方向总结}
\DeckKicker{方向总结}
\SlideMeta{证据等级}{方向综述}{来源状态：公开材料综合}
{\large\bfseries\color{Ink} Memory encryption / integrity / replay protection 的结论是：三篇主讲材料共同给出技术演进，但仍要用证据等级约束每个结论。\par}\vspace{1.8mm}
\begin{columns}[T,totalwidth=\textwidth]
  \begin{column}{0.45\textwidth}
    \fcolorbox{Rule}{Paper}{%
  \begin{minipage}[t]{0.97\linewidth}
    \vspace{1.1mm}
    \hspace{1.4mm}\begin{minipage}{0.92\linewidth}
      \PanelTitle{内容说明}
      \scriptsize \begin{itemize}
  \item Memory Encryption: A Survey of Existing Techniques：开创/基础入口，E2 / Background substrate。
  \item Bonsai Merkle Trees for memory integrity and freshness：代表性改进，E1 / Peer-reviewed mechanism anchor。
  \item AMD SEV-SNP industry memory-integrity evidence：当前 SOTA/补充边界，E4 / Industry evidence。
  \item 本方向结论：先用三篇主讲材料建立演进关系，再用证据等级控制机制、实验和商业化结论。
\end{itemize}
    \end{minipage}
    \vspace{1mm}
  \end{minipage}}
  \end{column}
  \begin{column}{0.49\textwidth}
    \fcolorbox{Rule}{Panel}{%
  \begin{minipage}[t]{0.97\linewidth}
    \vspace{1.1mm}
    \hspace{1.4mm}\begin{minipage}{0.92\linewidth}
      \PanelTitle{本方向方法对比}
      \scriptsize {\tiny
\begin{tabularx}{\linewidth}{@{}YYY@{}}
\toprule
\textbf{论文} & \textbf{定位} & \textbf{选择理由} \\
\midrule
Memory Encryption: A Survey of Existing Techniques & 开创/基础入口 & Henson 2014 is the background starting point for secure memory encryption。 \\
Bonsai Merkle Trees for memory integrity and freshness & 代表性改进 & Peer-reviewed Bonsai Merkle Tree mechanism anchor for memory integrity, f。 \\
AMD SEV-SNP industry memory-integrity evidence & 当前 SOTA/补充边界 & AMD SEV-SNP remains industry evidence for modern VM memory-integrity depl。 \\
\bottomrule
\end{tabularx}
}
    \end{minipage}
    \vspace{1mm}
  \end{minipage}}
  \end{column}
\end{columns}
\vspace{1mm}
{\tiny\textcolor{Muted}{引用依据：总结页综合三篇主讲材料的选择理由、评价字段和证据等级。}}
\end{frame}
